\section{讨论}
\label{sec:discussion}

\textbf{四次反转的共同结构。}每一个都是分母从未被建立的比值：没有零模型的对照、没有阳性预测率的召回率、分母未知的分数、跨两种输出体制的比较。修复是程序性的，而不是架构性的。四项审计都可以在 CPU 上、对基准本来就已写出的工件运行，只需几分钟。

\textbf{一份报告清单。}我们建议，任何形如``模型 $M$ 依赖图像''、``模型 $M$ 鲁棒''或``$M$ 可迁移''的主张都应附带：(1) 真实条件下的输出支撑大小与众数占比；(2) 对照的真实图像间零模型；(3) 阳性召回率、阳性预测率与精确率；(4) 契约覆盖与输出模式指示量；(5) 有效输入分辨率与精度；(6) 每个指标的分母。前四项是前提，后两项是测量声明。

\textbf{为什么文献需要重读。}此前的无图像消融报告对照，却没有零模型，也没有输出支撑检查。我们的结果表明，这些数字至少有两种误导方式：拒答模板产生很大的对照却没有证据提取；坍缩模型因为缺乏变化而产生零对照。这不是对任何具体论文执行的批评；而是在说，按常见方式计算的统计量并不能识别目标量。

\textbf{静默失效是与部署最相关的象限。}如果监控系统监视不确定性、拒答或推理通道激活，它能抓住分辨率损失与模糊，却会漏掉噪声。这种不对称不是某个模型的性质：五个模型的合并标记率在分辨率下为 14.40\%，在噪声下为 0.093\%，而召回损伤相当。我们谨慎界定它说明与不说明什么。标记率是\emph{模型侧}信号，不是人类或图像侧可见性的度量；在同一批数据上，一个无参考 Laplacian 方差统计量以 AUC $1.000$ 区分干净与重度噪声，基于规则的搜索在任何重度噪声回答中都找不到显式质量语言（附录~\ref{app:risk}，表~\ref{tab:imageside}，以及 \url{scripts/83_quality_language_audit.py}）。因此，主张更窄，也更有用：\emph{模型自身的监控与报告信号检测不到一种已被证明会损害其答案的退化}。实际监控器应把图像侧退化检测器与模型侧犹豫信号结合起来，因为两者在退化谱的两端失效。

\textbf{契约是模型可观测行为的一部分。}模式消融表明，同一组权重可以因所请求的输出契约而成为``自信正常''的捷径，或成为弃答模型。这既是测量警告（契约是一个变量），也是设计杠杆（若希望弃答，可以请求弃答）。它还意味着``模型对空白图像过度自信''不是一个模型级事实。

\textbf{我们未能建立的东西。}激发本研究的 selective-trust 框架未获验证：其组件只使错误排序发生边际变化，且不产生覆盖率。我们把它作为负结果报告，而不是在看到测试集后拟合一个更宽松的阈值。
